\EnglishSourceFrame{GELFAND-P001359}{unlabelled}
\underline{23.3.87}\qquad \underline{Gelfand seminar.}

\underline{\underline{V. E. Zakharov.}}\qquad \underline{$\bar{\partial}$-problems and integrable systems.}

Construction of integrable equations by the inverse-problem method.
(actually neither integrable nor by the inverse-problem method)

1973\quad A. B. Shabat formulated the dressing method.

Nonlinear equations are considered as compatibility conditions for overdetermined linear systems
\begin{equation*}
u_t+6uu_x+u_{xxx}=0\qquad (\text{KdV}) \tag{1}
\end{equation*}
--- the compatibility condition for the equations
\begin{equation*}
\begin{cases}
\psi_{xx}+u\psi=k^2\psi,\\
\psi_t=4\psi_{xxx}+6u\psi_x+3u_x\psi.
\end{cases} \tag{2}
\end{equation*}
\EnglishUnclear{GELFAND-P001359}{eq2-signs}{The signs in (1), (2) and the exponent in (7) are retained as written in the manuscript; their consistency requires checking.}

The method consists in considering the equation
\begin{equation*}
F(x+z,t)+\int K(x,s,t)F(s+z,t)\,ds+K(x,z,t)=0. \tag{3}
\end{equation*}
[In the arguments of $K$, plus signs $+$ are crossed out, with commas placed beneath them. The letters $K$ and $F$ under the integral are written over earlier writing:
\EnglishUnclear{GELFAND-P001359}{eq3-overwriting}{the original letters are illegible}.]

for any function $F$ for which
\begin{equation*}
\frac{\partial}{\partial t}F(\xi,t)=8\frac{\partial^3F}{\partial\xi^3}. \tag{4}
\end{equation*}
In (3), the integral is understood in one of three ways:
\begin{equation*}
\int=\int_x^\infty\quad\text{or}\quad\int_{-\infty}^x
\quad\text{or}\quad
\frac12\left(\int_x^\infty\ldots-\int_{-\infty}^x\ldots\right). \tag{5}
\end{equation*}
Then
\begin{equation*}
u=2\frac{d}{dx}K(x,x,t)\quad\text{--- a solution of (1)}. \tag{6}
\end{equation*}
Taking
\begin{equation*}
\psi_0(x,t)=e^{ikx+4ik^3t}, \tag{7}
\end{equation*}
and then
\begin{equation*}
\psi=\psi_0(x)+\int K(x,s)\psi_0(s)\,ds, \tag{8}
\end{equation*}
we obtain [crossed out: \EnglishUnclear{GELFAND-P001359}{after-eq8-cancelled}{illegible, possibly $\psi$}]
--- a solution of [inserted above the line: the system] (2), from which $u$ is found.

Consider also more general functions $F$ satisfying the equations
\begin{equation*}
\begin{cases}
\dfrac{\partial^2F}{\partial x^2}-\dfrac{\partial^2F}{\partial z^2}=0,\\[6pt]
\dfrac{\partial F}{\partial t}=4\left(\dfrac{\partial^3F}{\partial x^3}+\dfrac{\partial^3F}{\partial z^3}\right).
\end{cases} \tag{9}
\end{equation*}


\par\smallskip\noindent\textit{Reading notes.}\par
\noindent\texttt{GELFAND-0148-U001} (eq2-signs): In (1), the signs before 6uu\_x and u\_xxx are plus signs; (2) reads psi\_xx+u psi=k\textasciicircum{}2 psi and psi\_t=4 psi\_xxx+6u psi\_x+3u\_x psi. In (7), region [400,2320,1160,2480], the exponent is ikx+4ik\textasciicircum{}3t.\par
Candidates: The literal writing of (1), (2) and (7), with the visible plus signs. Unresolved in the diplomatic source.\par
For u=0, the function in (7) gives psi\_xx=-k\textasciicircum{}2 psi and psi\_t=+4ik\textasciicircum{}3 psi, whereas 4 psi\_xxx=-4ik\textasciicircum{}3 psi. This is an internal mathematical check, not a correction to the source. The sign contradiction is retained. \par
\noindent\texttt{GELFAND-0148-U002} (eq3-overwriting): In (3), the first K and following F are retraced or corrected with dense strokes. Between x and s, and between x and z, crossed-out plus signs have commas beneath them.\par
Candidates: Originally F beneath the final K and K beneath the final F; The original letters are illegible. Unresolved in the diplomatic source.\par
The final writing K(x,s,t)F(s+z,t) and K(x,z,t) is readable; the underlying letters have not been recovered. An interchange of letters is not asserted.\par
\noindent\texttt{GELFAND-0148-U003} (after-eq8-cancelled): After "we obtain" and before the dash, there is short densely crossed-out writing.\par
Candidates: psi; an illegible crossed-out symbol. Unresolved in the diplomatic source.\par
The individual strokes do not permit confident recovery of the cancelled writing.\par


\EnglishSourceFrame{GELFAND-P001360}{24}
Then everything still holds as before.

\underline{Generalization.} [Wavy underlining.] In place of $(9)_1$, write:
\begin{equation*}
\frac{\partial^2F}{\partial x^2}-\frac{\partial^2F}{\partial z^2}
=\alpha\frac{\partial F}{\partial y}. \tag{10}
\end{equation*}
Then
\begin{equation*}
\begin{cases}
\psi_{xx}+u\psi=\alpha\dfrac{\partial\psi}{\partial y},\\[4pt]
\psi_t=4\psi_{xxx}+6u\psi_x+3u_x\psi
\end{cases} \tag{11}
\end{equation*}
\EnglishUnclear{GELFAND-P001360}{eq11-coefficient}{The coefficient of $\partial\psi/\partial y$: $\alpha$ or $2$.}

and for $u$ we obtain:
\begin{equation*}
\frac{\partial}{\partial x}\left[u_t+6uu_x+u_{xxx}\right]=3\alpha^2u_{yy}. \tag{12}
\end{equation*}
$y$ becomes a new parameter in (3).

This was in the 1974 paper of Zakharov--Shabat.

The third integral in (5) is defined when the function decays in both directions.
However, one may [crossed out:
\EnglishUnclear{GELFAND-P001360}{before-eq13-short-cancelled}{an illegible beginning of a word}]
understand the integral differently:

[Crossed out on the line below:
\EnglishUnclear{GELFAND-P001360}{before-eq13-words-cancelled}{illegible; possibly "as Fourier"}.]
\begin{equation*}
\int e^{\lambda x}\,dx=\frac1\lambda e^{\lambda x}\qquad\text{and so on} \tag{13}
\end{equation*}
We shall understand the integral in the third way in (5).
Introduce the function $T$ by the formula
\begin{equation*}
F(x,z)=\frac1{2\pi}\iint
T(\lambda',\bar\lambda',\lambda,\bar\lambda,t,y)
 e^{i(\lambda'x-\lambda z)}\,d\lambda'\,d\bar\lambda'\,d\lambda\,d\bar\lambda. \tag{14}
\end{equation*}
[Crossed out after $\bar\lambda'$:
\EnglishUnclear{GELFAND-P001360}{eq14-cancelled-argument}{an illegible symbol}.
The arguments $t,y$ are inserted above the line before the closing parenthesis.]

and define the function $K$ as
\begin{equation*}
K(x,z)=-\frac1{2\pi}\int K(\lambda,\bar\lambda,x,t,y)
 e^{i\lambda(x-z)}\,d\lambda\,d\bar\lambda. \tag{15}
\end{equation*}
[On the left-hand side, the second argument is overwritten:
\EnglishUnclear{GELFAND-P001360}{eq15-overwritten-argument}{the final $z$; the earlier writing is illegible}.
The arguments $t,y$ on the right are inserted above the line.]

This is a way to find growing solutions of the G.-L.-M. equation.

If $K$ and $F$ satisfy the integral equation, everything is in order
(using convention (13)).

[Crossed out: $\chi(x,y)$.] Set:
\begin{equation*}
\chi(\lambda,\bar\lambda,x,y,t)=1+\frac1{2\pi i}
\int\frac{K(\eta,\bar\eta,x,y,t)}{\lambda-\eta}\,d\eta\,d\bar\eta. \tag{16}
\end{equation*}
Here:
\begin{equation*}
\frac1{\lambda-\eta}=\lim_{\varepsilon\to0}
\frac{\bar\lambda-\bar\eta}{|\lambda-\eta|^2+\varepsilon^2}. \tag{17}
\end{equation*}
Write the equation:


\par\smallskip\noindent\textit{Reading notes.}\par
\noindent\texttt{GELFAND-0148-U004} (eq11-coefficient): Before the y derivative of psi in the first line of (11), a sign has a small upper loop and crossing lower strokes; it resembles both alpha and 2.\par
Candidates: alpha; 2. Unresolved in the diplomatic source.\par
Alpha is readable in (10) and (12), making alpha contextually likely without resolving the graphical ambiguity in (11). The formula displays the candidate alpha with an explicit marker; this is not a correction.\par
\noindent\texttt{GELFAND-0148-U005} (before-eq13-short-cancelled): Between "differently" and "understand", a short beginning of a word is crossed out.\par
Candidates: ob; op; illegible. Unresolved in the diplomatic source.\par
The earlier writing has not been recovered.\par
\noindent\texttt{GELFAND-0148-U006} (before-eq13-words-cancelled): Beneath the line beginning "However, one may...", two short words are crossed out. The second has a loop resembling the Cyrillic letter Ф and a descender.\par
Candidates: as Fourier; by Fourier; two illegible words. Unresolved in the diplomatic source.\par
Both verbal candidates are conjectural; the crossed-out text has not been silently omitted.\par
\noindent\texttt{GELFAND-0148-U007} (eq14-cancelled-argument): In (14), between the arguments lambda-bar-prime and lambda, a cancelled symbol is crossed by several slanted strokes.\par
Candidates: lambda; chi; an illegible symbol. Unresolved in the diplomatic source.\par
Analysis of the arguments of T does not recover the cancelled symbol; it is recorded separately from the final formula.\par
\noindent\texttt{GELFAND-0148-U008} (eq15-overwritten-argument): On the left-hand side of (15), the second argument of K is densely retraced. The final z is visible, with additional strokes beneath it.\par
Candidates: z over x; a retraced z. Unresolved in the diplomatic source.\par
The final z is preserved; the original writing is not asserted.\par


\EnglishSourceFrame{GELFAND-P001361}{unlabelled}
\begin{equation*}
\frac{\partial\chi}{\partial\bar\lambda}
=\int\chi(\eta,\bar\eta,x_i)T(\eta,\bar\eta,\lambda,\bar\lambda,x_i)\,d\eta\,d\bar\eta. \tag{18}
\end{equation*}
\begin{equation*}
x_i=(x,y,t)
\end{equation*}
This is equivalent to the original approach.

\underline{The dressing method in the language of $\bar\partial$-problems.}
\EnglishUnclear{GELFAND-P001361}{heading-ending}{The heading ending reads "zadach." or "zadachi." [different Russian endings for "problem(s)"].}

Let $T$ satisfy the system.
\begin{equation*}
\begin{cases}
\dfrac{\partial T}{\partial x}=-i(\lambda-\eta)T,\\[5pt]
\dfrac{\partial T}{\partial y}=-\alpha(\lambda^2-\eta^2)T,\\[5pt]
\dfrac{\partial T}{\partial t}=-4(\lambda^3-\eta^3)T,
\end{cases} \tag{19}
\end{equation*}
from which
\begin{equation*}
T=T_0(\eta,\bar\eta,\lambda,\bar\lambda)
 e^{-i(\lambda-\eta)x-\alpha(\lambda^2-\eta^2)y-4(\lambda^3-\eta^3)t}. \tag{20}
\end{equation*}
In place of (19), one may introduce $x=x_i=(x_1,x_2,x_3)$ and the functions

$K_i(\lambda)$ and $I_i(\lambda)$:

[Crossed out in (21):
\EnglishUnclear{GELFAND-P001361}{eq21-cancelled}{possibly $[K_i(\lambda),K_i(\eta)]=0$; the subscripts are uncertain}.]
\begin{equation*}
[I_i(\lambda),I_j(\eta)]=0. \tag{21}
\end{equation*}
\begin{equation*}
\frac{\partial T}{\partial x_i}=I_i(\eta)T-TI_i(\lambda). \tag{22}
\end{equation*}
In the example
\begin{equation*}
I_1=i\lambda,\quad I_2=\alpha\lambda^2,\quad I_3=4\lambda^3. \tag{23}
\end{equation*}
\EnglishUnclear{GELFAND-P001361}{time-coefficient}{In (19), (20) and (23), the coefficient of the cubic powers is written as $4$, without $i$; this is preserved.}

Now solve the nonlocal $\bar\partial$-problem (18), with the normalization condition $\chi(\infty)=1$.

Let
\begin{equation*}
\chi=1+\frac{\chi_0(x,y,t)}{\lambda}+\ldots \tag{24}
\end{equation*}
[Crossed out in the numerator, before $x,y,t$:
\EnglishUnclear{GELFAND-P001361}{eq24-cancelled-argument}{an illegible argument, possibly $\lambda$}.]

Then
\begin{equation*}
u=2i\frac{\partial\chi_0}{\partial x} \tag{25}
\end{equation*}
satisfies [crossed out:
\EnglishUnclear{GELFAND-P001361}{eq25-cancelled-word}{possibly "formula" [source abbreviation]}]
the KP equation: (12).


\par\smallskip\noindent\textit{Reading notes.}\par
\noindent\texttt{GELFAND-0148-U009} (heading-ending): The last word of the underlined heading after d-bar has an ending in which the last stroke of ch is difficult to distinguish from an additional i.\par
Candidates: zadach.; zadachi. Unresolved in the diplomatic source.\par
The grammatical preference for the plural is not treated as resolution of the graphical doubt.\par
\noindent\texttt{GELFAND-0148-U010} (eq21-cancelled): To the left of the final commutator (21), a commutator of two K entries with arguments lambda and eta, an equality sign and a final digit are crossed out. Strokes cross the subscripts.\par
Candidates: [K\_i(lambda),K\_i(eta)]=0; [K\_i(lambda),K\_j(eta)]=0. Unresolved in the diplomatic source.\par
In both K entries the clearest subscript resembles i; replacement of the second by j by analogy with the final formula is not asserted.\par
\noindent\texttt{GELFAND-0148-U011} (time-coefficient): The third line of (19) visibly has the coefficient -4 without i. The same occurs in the exponent (20), region [1510,1310,1810,1450], and in I\_3=4 lambda\textasciicircum{}3 in (23), region [1135,2400,1420,2530].\par
Candidates: -4(lambda\textasciicircum{}3-eta\textasciicircum{}3) in (19),(20); I\_3=4 lambda\textasciicircum{}3 in (23). Unresolved in the diplomatic source.\par
Formulas (19),(20),(22),(23) agree with one another, but the connection of the time normalization to (7),(14) and the claimed KP equation (12) requires separate checking. No factor i was added from mathematical expectation.\par
\noindent\texttt{GELFAND-0148-U012} (eq24-cancelled-argument): After the opening parenthesis of chi\_0 and before x,y,t, an additional argument is densely crossed out.\par
Candidates: lambda; eta; an illegible argument. Unresolved in the diplomatic source.\par
The remaining arguments x,y,t are readable; the cancelled sign has not been recovered.\par
\noindent\texttt{GELFAND-0148-U013} (eq25-cancelled-word): Between "satisfies" and "the equation", a short word or abbreviation is crossed out; its first letter has a descender.\par
Candidates: formula [abbreviated]; function [abbreviated]; illegible. Unresolved in the diplomatic source.\par
The conjectural "formula" abbreviation is not asserted as an established reading.\par


\EnglishSourceFrame{GELFAND-P001362}{25}
(26)\quad \underline{Main assertion:} [Wavy underlining.]

For $K_j$, one may take any three commuting rational matrix-valued functions.

One may also take functions on a Riemann surface.

(27)\quad \underline{New proposal.} [Wavy underlining.]

take $K_i(\lambda)$ to be commuting differential operators with coefficients rational in $\lambda$.
\begin{equation*}
L_i=\frac{\partial}{\partial x_i}+K_i(\lambda)+f_i(\lambda)\frac{\partial}{\partial\lambda}. \tag{28}
\end{equation*}
[A bracket beneath $f_i(\lambda)$ has the caption: scalar functions, rational, with freely moving \EnglishUnclear{GELFAND-P001362}{eq28-caption-ending}{poles [genitive ending] / poles [instrumental ending]}.]

(29)\quad $[L_i,L_j]=0\ \Rightarrow$ conditions on $K_i$ in which the number of equations =
\newline
= the number of unknowns.

Many of the resulting equations are new (if everything has variable coefficients)

The same can be done on an elliptic curve.

In the work of Belinsky/Zakharov (1978), a special case is gravitation.

\underline{New this year:}
\begin{equation*}
L_i=\frac{\partial}{\partial x_i}+M_i\left(\lambda,\frac{\partial}{\partial\lambda}\right). \tag{30}
\end{equation*}
\begin{equation*}
[L_i,L_j]=0\quad\Longleftrightarrow\quad[M_i,M_j]=0. \tag{31}
\end{equation*}
(32)\quad \underline{Example:} [Wavy underlining.]
\begin{equation*}
\begin{gathered}
M_1=I\frac{\partial}{\partial\lambda}+\frac A\lambda,
\qquad I^2=1,\quad A^2=-1,\quad IA+AI=0;\\[4pt]
M_2=M_1\frac{\partial}{\partial\lambda}+\frac C{\lambda^2},
\qquad C=-\frac12(I+A).
\end{gathered}
\end{equation*}
[On the right, beyond a vertical line:]

($I,A$ are matrices, for example,
\begin{equation*}
I=\begin{pmatrix}1&0\\0&-1\end{pmatrix},\qquad
A=\begin{pmatrix}0&i\\i&0\end{pmatrix}
\end{equation*}
\EnglishUnclear{GELFAND-P001362}{eq32-matrix-a}{The off-diagonal entries of $A$: $i,i$ or $1,1$.})

\begin{equation*}
L_1=\frac{\partial}{\partial x_1}+M_1,\qquad
L_2=\frac{\partial}{\partial x_2}+M_1^2,\qquad
L_3=\frac{\partial}{\partial x_3}+M_2.
\end{equation*}
A large amount of calculation is needed.

Finding solutions of the equations:
\begin{equation*}
L_i y=0. \tag{33}
\end{equation*}


\par\smallskip\noindent\textit{Reading notes.}\par
\noindent\texttt{GELFAND-0148-U014} (eq28-caption-ending): The last word of the caption beneath f\_i(lambda), following "with freely moving", ends with dense strokes resembling "ov"; the beginning "polyus" (pole) is readable.\par
Candidates: poles [genitive ending]; poles [instrumental ending]. Unresolved in the diplomatic source.\par
The grammatically expected instrumental ending is not substituted for the possible genitive ending. The text ends before the image edge; this is not loss at the edge.\par
\noindent\texttt{GELFAND-0148-U015} (eq32-matrix-a): Both off-diagonal signs in A have a slanted upper part and a lower horizontal stroke; a separate dot on i is unclear and merges with the paper grid. A\textasciicircum{}2=-1 is explicitly written nearby.\par
Candidates: A=((0,i),(i,0)); A=((0,1),(1,0)). Unresolved in the diplomatic source.\par
The i candidate gives A\textasciicircum{}2=-1 and, with I=diag(1,-1), satisfies IA+AI=0; the 1 candidate gives A\textasciicircum{}2=+1. This check supports i without resolving the graphical doubt. The body explicitly marks i as a candidate. Under the abstract identities I\textasciicircum{}2=1, A\textasciicircum{}2=-1 and IA+AI=0, the example commutator of M\_1 and M\_2 vanishes; no corrected matrix is asserted.\par


\EnglishSourceFrame{GELFAND-P001363}{unlabelled}
Start with the G.-L.-M. integral equations in $\lambda$, take $n$-soliton solutions, and degenerate them into rational ones.

\begin{center}
\begin{tikzcd}
M_i(\lambda,\frac{\partial}{\partial\lambda},x_i)\arrow[d]\\
M_i^0(\frac{\partial}{\partial\lambda})
\end{tikzcd}
\end{center}
construct $K(\lambda,\lambda',x_i)$ and $F(\lambda,\lambda',x_i)$

and then \EnglishUnclear{GELFAND-P001363}{after-diagram-verb}{it remains / it remained}
to construct rational solutions.

Dissipative systems are also included here.

The next generalization:
\begin{equation*}
L_i=R_i\left(\lambda,\frac{\partial}{\partial\lambda}\right)
\frac{\partial}{\partial x_i}
-M_i\left(\lambda,\frac{\partial}{\partial\lambda}\right). \tag{34}
\end{equation*}
\EnglishUnclear{GELFAND-P001363}{eq34-number}{The number (34) has superimposed strokes.}

\underline{\underline{A. A. Migdal.}}

\underline{Another meaning of}
[crossed out: \EnglishUnclear{GELFAND-P001363}{title-cancelled}{an illegible short beginning}]
\underline{the evaluation of matrix integrals and the passage}
\newline
\underline{from matrix integrals to continuum integrals}

["matrix" is inserted above the line before "integrals".]

A noncommutative analogue of
\EnglishUnclear{GELFAND-P001363}{analog-phrase}{possibly Karasev's symbol}.

(Operators directly, in place of $N\times N$ matrices)
\begin{equation*}
\int dA\,\exp\left(-\frac N\lambda\operatorname{tr}V(A)\right),
\qquad A^+=-A,\quad N\to\infty. \tag{1}
\end{equation*}
[Callout from $V(A)$: a polynomial in $A$ (arbitrary); these integrals are evaluated exactly as $N\to\infty$.]

\underline{A way to evaluate them.} [Wavy underlining.] Consider
\begin{equation*}
\psi(t)=\frac1N\operatorname{tr}(\exp(At)) \tag{2}
\end{equation*}
and derive the equation of motion (Schwinger--Dyson). Clearly,
\begin{equation*}
0=\int dA\,\frac1{N^2}\operatorname{tr}\left(
\frac{\partial}{\partial A}\left[e^{At}
\exp\left(-\frac N\lambda\operatorname{tr}V(A)\right)\right]\right). \tag{3}
\end{equation*}
[Crossed out between $e^{At}$ and $\exp$:
\EnglishUnclear{GELFAND-P001363}{eq3-cancelled-brackets}{a bracket or several brackets}.]

(the integral of a total derivative $=0$). This is equivalent to
\begin{equation*}
0=\int_0^t d\tau\,\psi(\tau)\psi(t-\tau)
-\frac1\lambda V'\left(\frac{\partial}{\partial t}\right)\psi(t), \tag{4}
\end{equation*}
\begin{equation*}
R(\omega)=\psi^2(\omega)-\frac1\lambda V'(\omega)\psi(\omega). \tag{5}
\end{equation*}

The object of interest is the expectation $\langle f(A)\rangle$ defined by the integral, rather than the integral itself.
\begin{equation*}
\langle f(A)\rangle=
\frac{\int dA\,\exp(\ldots)f(A)}{\int dA\,\exp(\ldots)} \tag{6}
\end{equation*}
[On the right:]

($f(A)$ is some invariant function ---
\newline
--- a function of the eigenvalues


\par\smallskip\noindent\textit{Reading notes.}\par
\noindent\texttt{GELFAND-0148-U016} (after-diagram-verb): The word between "and then" and "construct" has a joined ending; its stem "osta..." is readable.\par
Candidates: it remains; it remained. Unresolved in the diplomatic source.\par
Neither a grammatical correction nor a confident reading of the ending was introduced.\par
\noindent\texttt{GELFAND-0148-U017} (eq34-number): The number (34) is readable to the left of the formula, with slanted and descending strokes over its digits.\par
Candidates: (34), retraced number; (34), crossed-out or corrected number. Unresolved in the diplomatic source.\par
The value 34 is preserved; the purpose of the additional strokes is not established.\par
\noindent\texttt{GELFAND-0148-U018} (title-cancelled): In Migdal's heading, after "Another meaning" and before "evaluation", a short beginning is crossed by several horizontal strokes.\par
Candidates: po; pod; illegible. Unresolved in the diplomatic source.\par
"Matrix" is inserted above the line; the final "continuum" was checked at 6x.\par
\noindent\texttt{GELFAND-0148-U019} (analog-phrase): After "A noncommutative analogue", there are two words. The first resembles "symbol"; the second begins with capital K and has the outline of "Karasev".\par
Candidates: Karasev's symbol; Karasev's meaning. Unresolved in the diplomatic source.\par
Context or external biographical information did not replace the reading. At final comparison, the surname was compared with the clear "Karasev's method" on P001366; uncertainty in the preceding word remains.\par
\noindent\texttt{GELFAND-0148-U020} (eq3-cancelled-brackets): After e\textasciicircum{}\{At\} in (3), a tall crossed-out sign resembles a closing bracket, possibly with an additional round parenthesis.\par
Candidates: ]; ]); illegible brackets. Unresolved in the diplomatic source.\par
The final outer round and square brackets are reproduced from the right-hand part of the line; the cancelled material is described separately.\par


\EnglishSourceFrame{GELFAND-P001364}{26}
I. M. G. \underline{Problem.}
(T. \EnglishUnclear{GELFAND-P001364}{top-name}{Alekseevskaya / Aleksievskaya})
(\EnglishUnclear{GELFAND-P001364}{top-parenthesis}{did [masculine] / did [feminine] / illegible})

Given a cube (3-dimensional). Pass a plane through each three of its vertices. Into how many parts do they divide the cube?

(\underline{Answer:} 96)

The limit as $N\to\infty$ is the same for all classical groups:
\begin{equation*}
O(\infty)=SU(\infty)=SP(\infty). \tag{7}
\end{equation*}
\begin{equation*}
0=\int dA\,\frac1{N^2}\frac{\partial}{\partial A_{ij}}
\left\{(e^{At})^{ji}\exp\left(-\frac N\lambda\operatorname{tr}V(A)\right)\right\}. \tag{3'}
\end{equation*}
[Crossed out after $1/N^2$:
\EnglishUnclear{GELFAND-P001364}{eq3-prime-cancelled}{possibly $\operatorname{tr}$}.]

The formula used was:
\begin{equation*}
\frac{\partial}{\partial A_{ij}}(e^{At})_{ji}
=\int_0^t d\tau\,(e^{A\tau})_{jj}(e^{A(t-\tau)})_{ii}. \tag{8}
\end{equation*}
Derivation: differentiate
\begin{equation*}
\prod(1+A\Delta t)=T\exp\left(\int A\,dt\right)
\end{equation*}
and obtain the integral (successively
\EnglishUnclear{GELFAND-P001364}{proof-verb}{pulling out [abbreviated] / illegible}
each matrix entry in its own place)

Another formula from matrix theory that was used:
\begin{equation*}
\frac{\partial}{\partial A_{ij}}\operatorname{Tr}V(A)=(V'(A))_{ji}. \tag{9}
\end{equation*}
From (8) and (9), (4) follows:
\begin{equation*}
\begin{aligned}
(e^{At})^{ji}\frac{\partial}{\partial A_{ij}}\operatorname{Tr}V(A)
&={}\\
&=(V'(A))_{ji}(e^{At})_{ij}
=V'\left(\frac{\partial}{\partial t}\right)\operatorname{tr}(e^{At}).
\end{aligned}
\end{equation*}
In (4), the expectation should have been written ($\int dA$ is omitted).

It was stated that
\begin{equation*}
0=\int_0^t d\tau\,\langle\psi(\tau)\psi(t-\tau)\rangle
-\frac1\lambda V'\left(\frac{\partial}{\partial t}\right)\langle\psi(t)\rangle. \tag{10}
\end{equation*}


\par\smallskip\noindent\textit{Reading notes.}\par
\noindent\texttt{GELFAND-0148-U021} (top-name): In parentheses after "I. M. G. Problem", the initial T. and a surname beginning "Aleks..." and ending "...vskaya" are written; two middle vowels are joined.\par
Candidates: T. Alekseevskaya; T. Aleksievskaya. Unresolved in the diplomatic source.\par
The surname was not recovered from information about the seminar; the source writing is retained as candidates.\par
\noindent\texttt{GELFAND-0148-U022} (top-parenthesis): After the surname, separate parentheses contain a short word with an unclear initial combination and an ending resembling "...lal" or "...lan".\par
Candidates: did [masculine]; did [feminine]; done; illegible. Unresolved in the diplomatic source.\par
This note belongs to the visible upper block and was not supplied from the catalogue. The answer 96 was checked separately at 8x in [600,305,755,435].\par
\noindent\texttt{GELFAND-0148-U023} (eq3-prime-cancelled): Between 1/N\textasciicircum{}2 and the partial derivative, a crossed-out cluster of two narrow signs may be an earlier tr.\par
Candidates: tr; illegible. Unresolved in the diplomatic source.\par
The upper ji on the exponential in (3-prime) and lower ji,jj,ii in (8) are reproduced in their respective positions, without standardizing the indices.\par
\noindent\texttt{GELFAND-0148-U024} (proof-verb): In the explanation of the product after "successively", an abbreviated word begins "vy..." and ends "...rg." or a similar combination.\par
Candidates: pulling out [abbreviated]; illegible. Unresolved in the diplomatic source.\par
The possible colloquial abbreviation for "pulling out" was not expanded to the complete Russian word or replaced by a mathematical explanation.\par


\EnglishSourceFrame{GELFAND-P001365}{unlabelled}
[Next, consider

["Next" is inserted above the crossed-out
\EnglishUnclear{GELFAND-P001365}{top-cancelled}{Now}.]
\begin{equation*}
\int\mathcal D A_\mu\,\exp\left(+\int d^4x\,\frac N\lambda
\operatorname{tr}F_{\mu\nu}^2\right) \tag{11}
\end{equation*}
]

Use the formula
\begin{equation*}
\left\langle\frac{\operatorname{tr}A}N\frac{\operatorname{tr}B}N\right\rangle
=\left\langle\frac{\operatorname{tr}A}N\right\rangle
\left\langle\frac{\operatorname{tr}B}N\right\rangle
+O\left(\frac1{N^2}\right). \tag{12}
\end{equation*}
Applying it, we obtain
\begin{equation*}
0=\int_0^t d\tau\,\langle\psi(\tau)\rangle\langle\psi(t-\tau)\rangle
-\frac1\lambda V'\left(\frac{\partial}{\partial t}\right)\langle\psi(t)\rangle. \tag{13}
\end{equation*}
Now introduce the Laplace transform:
\begin{equation*}
f(\omega)=\int_0^\infty dt\,e^{-\omega t}\langle\psi(t)\rangle; \tag{14}
\end{equation*}
then, from (13) (convolution!)
\newline
[crossed out: a] we obtain the equation
\begin{equation*}
\underbrace{R(\omega)}_{\text{a polynomial in }\omega}
=f^2(\omega)-\frac1\lambda V'(\omega)f(\omega) \tag{15}
\end{equation*}
or (taking jumps):
\begin{equation*}
0=2\operatorname{Im}f\left(\operatorname{Re}f-\frac1{2\lambda}V'(\omega)\right). \tag{16}
\end{equation*}
Here
\begin{equation*}
\operatorname{Im}f=\frac12\bigl(f(x+i0)-f(x-i0)\bigr). \tag{17}
\end{equation*}
But it is easy to see that
\begin{equation*}
f(\omega)=\frac1N\left\langle\operatorname{tr}\frac1{\omega-A}\right\rangle. \tag{18}
\end{equation*}
\begin{equation*}
f^2(x+i0)-f^2(x-i0)=2\operatorname{Im}f\cdot\operatorname{Re}f. \tag{19}
\end{equation*}
\EnglishUnclear{GELFAND-P001365}{jump-normalization}{In (17), $1/2$ is written; in (19), the factor is $2$. The notation is preserved.}

(20)\quad $f$
\EnglishUnclear{GELFAND-P001365}{eq20-description}{analytically continues into the left half-plane / is completely represented in the left half-plane}.

$f(\omega)$ has jumps at points in the spectrum of the operator $A$
(this is \EnglishUnclear{GELFAND-P001365}{last-parenthesis}{an infinite [abbreviated] / a real [abbreviated]} matrix)


\par\smallskip\noindent\textit{Reading notes.}\par
\noindent\texttt{GELFAND-0148-U025} (top-cancelled): Beneath the word "Next", inserted above the line, a crossed-out word has a tall initial stroke and a final descending stroke, resembling "Now".\par
Candidates: Now; illegible. Unresolved in the diplomatic source.\par
The left square bracket before "Next" and the closing bracket after (11) are reproduced. The plus sign in the exponent of (11) was checked and not corrected.\par
\noindent\texttt{GELFAND-0148-U026} (jump-normalization): In (17), Im f is defined as (1/2)(f(x+i0)-f(x-i0)); the denominator has no i. In (19), the difference of squares equals 2 Im f Re f.\par
Candidates: literal coefficients 1/2 and 2; a special unspecified normalization of Im and Re; a possible authorial inaccuracy in the coefficient or notation. Unresolved in the diplomatic source.\par
With the usual Re f=(f\_++f\_-)/2, definition (17) would give f\_+\textasciicircum{}2-f\_-\textasciicircum{}2=4 Im f Re f. Neither i nor the factor 4 was substituted in the text. The source may use another normalization, but none is explicitly defined on this page. Formula (16) is also retained literally.\par
\noindent\texttt{GELFAND-0148-U027} (eq20-description): In line (20), two closely written words or abbreviations follow f, then "into/in the left... half-plane...". The word for "left" has dense superimposed strokes.\par
Candidates: f analytically continues into the left half-plane; f is completely represented in the left half-plane. Unresolved in the diplomatic source.\par
Only f, the preposition "in/into" and the stem "half-plane" are securely distinguishable; candidates for the remaining phrase are not asserted. An expected property of the resolvent was not substituted.\par
\noindent\texttt{GELFAND-0148-U028} (last-parenthesis): In the final parentheses after "this is", a short abbreviation precedes "matrix"; a loop, an upper line and fine final strokes are visible.\par
Candidates: infinite [abbreviated]; real [abbreviated]. Unresolved in the diplomatic source.\par
Properties of A from preceding pages were not used to select a reading; contradictions were not removed.\par


\EnglishSourceFrame{GELFAND-P001366}{27}
Now the old (French) equation is obtained
\begin{equation*}
f(\omega)=\int\frac{\varrho(\lambda)\,d\lambda}{\omega-\lambda}. \tag{21}
\end{equation*}
\begin{equation*}
f=\frac1{2\lambda}V'(\omega)+M(\omega)\sqrt{(\omega-a)(\omega-b)}. \tag{22}
\end{equation*}
[\underline{Zakharevich.}
\begin{equation*}
A(x)(B(x)-\alpha(x))=0,\qquad
 iA(x)+B(x)\quad\text{analytically continues into }\operatorname{Im}x>0
\end{equation*}
and so on]

Instead, one may solve the quadratic equation (15):
\begin{equation*}
f=\frac1{2\lambda}V'(\omega)-\sqrt{\left(\frac{V'}{2\lambda}\right)^2-R}. \tag{23}
\end{equation*}
[Callout to $R$: the polynomial is chosen!]
\EnglishUnclear{GELFAND-P001366}{eq23-radicand}{The expression under the radical is written with $-R$.}
\begin{equation*}
f(\omega)\underset{\omega\to\infty}{\longrightarrow}
\frac1\omega-M(\omega)\sqrt{(\omega-a)(\omega-b)} \tag{24}
\end{equation*}
[Above the final expression: (from satisfaction of the boundary conditions).]
\EnglishUnclear{GELFAND-P001366}{eq24-asymptotic}{The right-hand side following the arrow is preserved in full, including $-M(\omega)\sqrt{(\omega-a)(\omega-b)}$.}

[Crossed-out line:
\EnglishUnclear{GELFAND-P001366}{crossed-heading}{Next, the problem / This equation / illegible}.]

For field theory, Karasev's method does not work:

[crossed out: \EnglishUnclear{GELFAND-P001366}{short-crossed}{ur [unexpanded abbreviation] / when}]
the components do not commute and cannot be brought simultaneously to diagonal form.

Consider the integral (11).
\begin{equation*}
A_\mu(x):\ \longrightarrow\ W(c)=\frac1N\operatorname{tr}
 P\exp\left(\oint_c A_\mu\,dx_\mu\right) \tag{25}
\end{equation*}
and repeat the same reasoning

$F_{\mu\nu}$ is a tensor in $\mu,\nu$ and a matrix in the Lie algebra of $SU(N)$.
\begin{equation*}
x_\mu=u_\mu(t),\qquad t\in(0,2\pi)
\end{equation*}
\EnglishUnclear{GELFAND-P001366}{curve-parameter}{The letter on the right of the equality: $u_\mu(t)$ or $x_\mu(t)$.}
\begin{equation*}
\int A_\mu(x(t))\dot x_\mu\,dt
\end{equation*}
\begin{equation*}
\sim\exp At\qquad\text{earlier}
\end{equation*}
A chain of identities: $\ldots$


\par\smallskip\noindent\textit{Reading notes.}\par
\noindent\texttt{GELFAND-0148-U029} (eq23-radicand): In (23), the sign before the radical is minus; its radicand is (V-prime/(2 lambda))\textasciicircum{}2-R. In the preceding frame, (15) reads R=f\textasciicircum{}2-V-prime f/lambda.\par
Candidates: literal reading of (23): minus R; a possible change in the sign convention for R or a source error. Unresolved in the diplomatic source.\par
The vertical stroke belongs to the grid. Quadratic equation (15), with its written definition of R, gives the radicand (V-prime/(2 lambda))\textasciicircum{}2+R. This plus was not inserted into the body. A nearby callout "the polynomial is chosen!" points to R.\par
\noindent\texttt{GELFAND-0148-U030} (eq24-asymptotic): After the arrow with omega tending to infinity beneath it, the consecutive expression is 1/omega-M(omega)sqrt((omega-a)(omega-b)). Above the second part is the caption "from satisfaction of the boundary conditions".\par
Candidates: the whole right-hand side belongs to the arrow; the condition f(omega) tends to 1/omega, followed by a separately added term containing M. Unresolved in the diplomatic source.\par
No visible separator lies between 1/omega and the minus, so the M term is retained in this formula rather than moved to (22). The meaning of the asymptotic condition remains unspecified.\par
\noindent\texttt{GELFAND-0148-U031} (crossed-heading): Before "For field theory", an entire two-word line is crossed out; the first word begins with D and the second has a low initial loop.\par
Candidates: Next, the problem; This equation; illegible. Unresolved in the diplomatic source.\par
These candidates are not a recovered heading; the cancelled source line is represented explicitly as crossed out.\par
\noindent\texttt{GELFAND-0148-U032} (short-crossed): Before "the components", a short crossed-out combination has a long lower stroke.\par
Candidates: ur [unexpanded abbreviation]; when; illegible. Unresolved in the diplomatic source.\par
The remaining text about noncommuting components, simultaneous reduction and diagonal form was read separately at 5x and 6x.\par
\noindent\texttt{GELFAND-0148-U033} (curve-parameter): In the parametrization x\_mu=...\_mu(t), the letter on the right consists of two thin slanted strokes resembling u; x(t) appears below in the argument of A.\par
Candidates: x\_mu=u\_mu(t); x\_mu=x\_mu(t). Unresolved in the diplomatic source.\par
The later x(t) was not used to replace the ambiguous letter. The lower parameter bound reads as an open 0, like other zeros in the manuscript.\par


\EnglishSourceFrame{GELFAND-P001367}{unlabelled}
\begin{equation*}
\int\mathcal D A\,\hat P\,\frac{\partial}{\partial A}
\exp\left(\oint_c A\,dy\right) \tag{26}
\end{equation*}
[Two short separator strokes.]

30 \underline{March} it is unknown whether the seminar will take place
(if cancelled, then
[crossed out: \EnglishUnclear{GELFAND-P001367}{announcement-cancelled}{s\ldots [initial Cyrillic s]}]
enquire at the laboratory, 139-53-57)


\par\smallskip\noindent\textit{Reading notes.}\par
\noindent\texttt{GELFAND-0148-U034} (announcement-cancelled): In the announcement, after "if cancelled, then" and before "enquire", a word is densely crossed out. Its initial Cyrillic s is readable at the left edge of the cancellation.\par
Candidates: s...; information; notify; illegible. Unresolved in the diplomatic source.\par
The rest of the announcement was read at 4x and 6x; the phone digits 139 were checked at 8x. The initial s is included in the cancelled material because the word boundary and cancellation do not coincide exactly. The rest of the frame is blank; no new date or speaker was added from catalogue associations.\par


\EnglishSourceFrame{GELFAND-P001368}{unlabelled}
\underline{$13/\overline{\mathrm{IV}}$-87}\quad \underline{Gelfand seminar}

\underline{\underline{Sinai}}

\underline{On spin glasses (following Parisi)}
[Wavy underlining.]
\begin{equation*}
H=\sum_{\langle i,j\rangle}J_{ij}\sigma_i\sigma_j;\qquad
J_{ij}\ \text{are random}. \tag{1}
\end{equation*}
(Ising model of a spin glass)

$\langle i,j\rangle$ is a pair of neighbours

\underline{Spin glass:} [Wavy underlining.]
when the interaction forces are long-range and strongly oscillatory

[Crossed out: \EnglishUnclear{GELFAND-P001368}{before-parizi-cancelled}{Par / Theor / illegible}.]
Parisi: all interact, namely:
\begin{equation*}
H=\frac1{\sqrt N}\sum J_{ij}\sigma_i\sigma_j \tag{2}
\end{equation*}
[On the right:] $J_{ij}$ is a Gaussian random variable
(all with the same variance)

There is a large number of \underline{ground states}
(differing in energy by $\sim1/N$);
["ground states" is underlined with a wavy line.]
\begin{equation*}
\sigma_i=\pm1.
\end{equation*}
Study the case $N\to\infty$ and $\displaystyle\lim_{N\to\infty}\min H$

[Beneath the limit, with a short connecting stroke: is calculated.]

\underline{Entropy of the set of configurations} $>0$ (in the limit $N\to\infty$)
[Inserted above the line between "set" and "configurations": (ground-state). Wavy underlining.]
\begin{equation*}
\frac{\ln\# [O(N)]}{N}\underset{N\to\infty}{\longrightarrow}h>0. \tag{3}
\end{equation*}
Property of \underline{ultrametricity} of ground states (Parisi, Virasoro): if
["ultrametricity" is underlined with a wavy line.]

$\sigma',\sigma'',\sigma'''$ are three ground states, then of the numbers
\begin{equation*}
\frac1N(\sigma',\sigma'');\quad
\frac1N(\sigma',\sigma''');\quad
\frac1N(\sigma'',\sigma''') \tag{4}
\end{equation*}
two must be equal, while the third is larger.

\par\smallskip\noindent\textit{Reading notes.}\par
\noindent\texttt{GELFAND-0148-U035} (before-parizi-cancelled): Immediately before "Parisi: all interact, namely", a short word is densely crossed out. A tall initial letter and several loops are visible, but the cancellation crosses all letters.\par
Candidates: Par; Theor; illegible. Unresolved in the diplomatic source.\par
The cancelled word is not established. The readable "Parisi" on the right was not copied into the crossed-out region by conjecture. Number (1) was checked at 8x and is securely 1; O(N) in (3) was checked at 6x and is an oval O without an inner stroke; the caption "is calculated" beneath the limit was checked at 5x. These latter places were not replaced by mathematically expected expressions.\par
